% ============================================================================
%  Actividad 5 · Puente de entrada de la parte de métricas y usabilidad
%
%  Este archivo abre la Actividad 5. NO introduce el documento completo: la
%  introducción general del volumen acumulado vive en main_completo.tex, y
%  duplicarla aquí dejaría dos aperturas compitiendo por el mismo papel.
%
%  Contiene tres piezas y ninguna más:
%    1. El puente desde la Actividad 4 y el orden interno de esta parte.
%    2. La nota de alcance de los resultados de usabilidad (participantes,
%       modalidad y por qué los dos bloques se analizan por separado).
%    3. \section{Descripción del producto} con \label{sec:a5-producto}, que es
%       el eco ejecutivo de la definición completa de la Actividad 1 y el
%       ancla del mapa de cumplimiento (contrato C4 de US-AVANCE-5).
% ============================================================================
\uxsection{Introducción}

Las cuatro entregas anteriores definieron el problema y la audiencia, tradujeron sus necesidades
en escenarios y recorridos, contrastaron la propuesta con las alternativas del mercado,
organizaron el contenido y construyeron un prototipo de alta fidelidad operable. Ninguna de ellas
midió. Esta quinta entrega cierra esa distancia: declara qué métrica de experiencia corresponde a
cada interfaz diseñada, aplica una prueba de usabilidad sobre el producto desplegado y convierte
lo observado en hallazgos ordenados y en un plan de corrección.

El orden va del marco a la evidencia y de la evidencia al cierre. El primer capítulo fija el marco
de medición: qué se mide en cada una de las siete interfaces, con qué instrumento, contra qué meta
y con qué resultado. El segundo documenta la prueba del jueves 20 de agosto de 2026 completa
---objetivos, participantes, escenarios, métricas, materiales, procedimiento, ejecución, análisis,
hallazgos y recomendaciones---. El tercero cierra el caso de estudio, consolida las referencias
bibliográficas y reúne los instrumentos de aplicación en el Anexo A. El Anexo B añade el manual de
uso del prototipo.

El alcance de la prueba comprende seis interfaces operables: acceso por perfil, búsqueda en el
catálogo, exportación, tablero ejecutivo, asistente conversacional y administración. Inicio y
gobierno del dato se recorrieron dentro de esas tareas y de los recorridos de trabajo, pero no
tuvieron tarea cronometrada propia; el capítulo de métricas lo declara fila por fila y reporta lo
que sí se observó en cada caso.

\begin{uxnota}[Alcance de los resultados de usabilidad]
La prueba reunió a cinco participantes reales. Carlos René Flores M. y Verónica Itzel Flores
González autorizaron aparecer por nombre y realizaron seis tareas estandarizadas. Otros tres
profesionales realizaron recorridos vinculados con su trabajo y solicitaron anonimato, por lo que
se identifican como A1, A2 y A3. Estas tres sesiones ocurrieron el 20 de agosto de 2026, una de
forma presencial y dos remotas, en computadora y con la misma versión del prototipo. Los dos
bloques se analizan por separado cuando sus métricas no son equivalentes.
\end{uxnota}

% ============================================================================
\section{Descripción del producto}\label{sec:a5-producto}

Este apartado resume el producto que la prueba midió. Su definición completa ---tipo de producto,
las doce características principales, los beneficios de experiencia, los principios de diseño, el
posicionamiento frente a las alternativas y los criterios de éxito--- es la de la Actividad 1 y no
se repite aquí.

Karisma Data responde a una dificultad institucional concreta: las fuentes de crédito, liquidez,
derivados y otros dominios se consultan en sistemas separados, con definiciones, permisos y
formatos que no siempre son comparables. Esto retrasa la validación de cifras, dificulta descubrir
fuentes relacionadas y hace que parte del conocimiento permanezca distribuido entre áreas.

La solución es un portal web modular que ofrece una entrada común y adapta la profundidad a la
tarea. Una consulta rápida empieza con el buscador; el análisis continúa con filtros, metadatos y
exportaciones; la supervisión utiliza indicadores consolidados; y las tareas de gobierno o
administración aparecen solo para los perfiles autorizados.

\begin{uxtabla}{L{3.5cm}Y}{Capacidades principales y valor para el usuario}
  \uxheadrow \thd{Capacidad} & \thd{Valor UX} \\
  Buscador y catálogo de datos & Localiza campos por concepto o nombre físico y muestra dominio, fuente, responsable y certificación \\
  Exploración y exportación & Permite continuar desde el hallazgo hasta un archivo reutilizable, con filtros y seguimiento del proceso \\
  Tableros e indicadores & Resume el último corte y permite profundizar en series y proyecciones \\
  Asistente conversacional & Responde en lenguaje natural y hace visible la consulta y la fuente que respaldan cada cifra \\
  Acceso por perfiles & Ajusta navegación, contenido y acciones sin crear productos separados para cada rol \\
  Gobierno y administración & Centraliza definiciones, permisos, estados de cuenta y evidencia de cambios \\
\end{uxtabla}

El valor esperado se resume en tres resultados: reducir el tiempo necesario para localizar y
validar información, conservar la procedencia cuando el dato se reutiliza y permitir que una misma
persona cambie de una vista simple a una herramienta profunda sin abandonar el portal.
